\section{Methodology}
\label{sec:method}

\subsection{Problem Formulation}

Let $\mathcal{D}=\{x_i\}_{i=1}^{N}$ denote the unlabeled distillation corpus. A teacher encoder $T$ maps each input to a cached teacher embedding
\begin{equation}
    t_i = T(x_i) \in \mathbb{R}^{d_T},
\end{equation}
and a compact student encoder $S_{\theta}$ maps the same input to
\begin{equation}
    s_i = S_{\theta}(x_i) \in \mathbb{R}^{d_S}.
\end{equation}
The objective is to train $S_{\theta}$ so that the student preserves the semantic organization induced by the teacher over the corpus, including nearest-neighbor rankings, local semantic regions, and higher-order community structure.

Standard pointwise embedding distillation minimizes a distance between each projected student embedding and its teacher embedding:
\begin{equation}
    \mathcal{L}_{\mathrm{point}}
    =
    \frac{1}{N}
    \sum_{i=1}^{N}
    \left(1-\cos(Ws_i,t_i)\right),
\end{equation}
where $W$ is an optional projection. Relational distillation instead matches student and teacher similarity matrices inside a mini-batch $B$:
\begin{equation}
    \mathcal{L}_{\mathrm{rel}}
    =
    \left\|
    G_B^S - G_B^T
    \right\|_F^2,
\end{equation}
where $G_B^S$ and $G_B^T$ are cosine-similarity matrices over examples in $B$. These objectives are useful local surrogates, but they do not directly supervise the corpus-level retrieval operator induced by the teacher.

\subsection{Motivation}

HeatGeo starts from the observation that retrieval quality depends on how probability mass is organized over an indexed corpus, not only on the absolute coordinate of each embedding. A compact student may be unable to reproduce the teacher's full embedding map, yet it can still be a strong retrieval model if it preserves the teacher's neighborhood distributions and global semantic topology.

We therefore treat the teacher embedding space as a semantic graph and distill the random-walk behavior of this graph. For an anchor $x_i$, a one-hop walk captures direct semantic neighbors; multi-hop walks reveal broader regions connected through intermediate examples. This makes the distillation target a soft corpus-level object: how semantic mass should spread from an anchor to nearby candidates at different scales.

For a diffusion scale $r$, write the full teacher diffusion row as
\begin{equation}
    q_{i,r}(j) = (P_r^T)_{ij}, \qquad j\in\mathcal{D}.
\end{equation}
Since matching a distribution over all $N$ examples is expensive, HeatGeo samples a candidate set $C_i$ and uses the conditional diffusion distribution on that set. Its captured mass is
\begin{equation}
    \rho_{i,r}(C_i)=\sum_{j\in C_i}q_{i,r}(j).
\end{equation}
The method is designed so that $C_i$ contains high-diffusion neighbors, hard negatives, and random negatives, making the restricted target informative while keeping training tractable.

The rest of the section defines the two geometric transfer losses and the stabilizing anchor objective. Formal proofs for the theoretical statements are provided in Appendix~\ref{app:proofs}.

\subsection{Multi-Scale Diffusion Loss}

HeatGeo first encodes the full training corpus with the teacher and constructs a teacher semantic graph. For each example $x_i$, let
\begin{equation}
    \mathcal{N}_k^T(i)
    =
    \operatorname{TopK}_{j\neq i}\cos(t_i,t_j)
\end{equation}
be its top-$k$ teacher nearest neighbors. To reduce noisy one-directional edges and hubness, HeatGeo uses a mutual $k$NN graph:
\begin{equation}
    j \in \mathcal{M}_k^T(i)
    \Longleftrightarrow
    j \in \mathcal{N}_k^T(i)
    \;\mathrm{and}\;
    i \in \mathcal{N}_k^T(j).
\end{equation}
The weighted adjacency matrix is
\begin{equation}
    W_{ij}^{T}
    =
    \begin{cases}
    \exp(\cos(t_i,t_j)/\tau_g), & j \in \mathcal{M}_k^T(i),\\
    0, & \mathrm{otherwise},
    \end{cases}
\end{equation}
where $\tau_g$ controls graph sharpness. Let $D^T$ be the diagonal degree matrix. The row-normalized transition matrix is
\begin{equation}
    P^T = (D^T)^{-1}W^T.
\end{equation}
HeatGeo then defines multi-scale diffusion targets
\begin{equation}
    P_r^T = (P^T)^r,
    \qquad r \in \mathcal{R},
\end{equation}
where $\mathcal{R}$ is typically $\{1,2,4\}$. Small $r$ preserves sharp local neighborhoods, while larger $r$ captures semantic regions and multi-hop communities.

For each anchor $x_i$, HeatGeo builds a candidate set
\begin{equation}
    C_i = C_i^{+} \cup C_i^{h} \cup C_i^{r},
\end{equation}
where $C_i^{+}$ contains high-probability diffusion neighbors, $C_i^{h}$ contains teacher hard negatives, and $C_i^{r}$ contains random negatives. The teacher diffusion target at scale $r$ is the diffusion mass restricted and renormalized over $C_i$:
\begin{equation}
    p_{i,r}^{T}(j)
    =
    \frac{(P_r^T)_{ij}}
    {\sum_{j' \in C_i}(P_r^T)_{ij'}},
    \qquad j\in C_i.
\end{equation}

Given an anchor embedding $s_i$ and candidate embeddings $\{s_j:j\in C_i\}$, the student induces a distribution over the same candidate set:
\begin{equation}
    p_i^{S}(j)
    =
    \frac{\exp(\cos(s_i,s_j)/\tau_s)}
    {\sum_{j'\in C_i}\exp(\cos(s_i,s_{j'})/\tau_s)},
    \qquad j\in C_i,
\end{equation}
where $\tau_s$ is the student temperature. The diffusion loss is
\begin{equation}
    \mathcal{L}_{\mathrm{diff}}
    =
    \frac{1}{|B|}
    \sum_{i\in B}
    \sum_{r\in\mathcal{R}_{e}}
    \omega_r
    D_{\mathrm{KL}}
    \left(
        p_{i,r}^{T} \| p_i^{S}
    \right),
\end{equation}
where $\omega_r$ is the scale weight and $\mathcal{R}_{e}$ is the set of active scales at epoch $e$.

\smallskip
\noindent\textbf{Lemma 1 (Diffusion KL preserves neighborhood mass).}
Fix an anchor $i$, a diffusion scale $r$, and candidate set $C_i$. If
\begin{equation}
    D_{\mathrm{KL}}\left(p^T_{i,r}\|p_i^S\right)\le \delta,
\end{equation}
then, for every subset $A\subseteq C_i$,
\begin{equation}
    \left|p_i^S(A)-p^T_{i,r}(A)\right|
    \le
    \sqrt{\frac{\delta}{2}}.
\end{equation}

\paragraph{Implication for HeatGeo.}
The subset $A$ can be any semantic slice of the candidate neighborhood: high-diffusion positives, teacher hard negatives, random negatives, a local cluster inside $C_i$, or a retrieval prefix induced by a downstream ranking rule. Thus $\mathcal{L}_{\mathrm{diff}}$ preserves how much semantic probability mass the teacher assigns to regions of the retrieval neighborhood, rather than merely aligning isolated vectors.

\smallskip
\noindent\textbf{Lemma 2 (Candidate restriction under high captured mass).}
Let $q_{i,r}(j)=(P_r^T)_{ij}$ be the full-corpus teacher diffusion row, let $\rho=\rho_{i,r}(C_i)>0$, and let $p^T_{i,r}$ be the restricted target on $C_i$. Define the zero extension of the student candidate distribution to the corpus by
\begin{equation}
    \widehat p_i^S(j)=
    \begin{cases}
    p_i^S(j), & j\in C_i,\\
    0, & j\notin C_i.
    \end{cases}
\end{equation}
Then
\begin{equation}
    \mathrm{TV}_{\mathcal{D}}\left(q_{i,r},\widehat p_i^S\right)
    \le
    (1-\rho)
    +\rho\,\mathrm{TV}_{C_i}\left(p^T_{i,r},p_i^S\right).
\end{equation}
Consequently, if $\rho_{i,r}(C_i)\ge 1-\epsilon$ and
$D_{\mathrm{KL}}\left(p^T_{i,r}\|p_i^S\right)\le \delta$, then
\begin{equation}
    \mathrm{TV}_{\mathcal{D}}\left(q_{i,r},\widehat p_i^S\right)
    \le
    \epsilon + \sqrt{\frac{\delta}{2}}.
\end{equation}

\paragraph{Implication for HeatGeo.}
The restricted target $p^T_{i,r}$ is the teacher diffusion row conditioned on the sampled candidates. Lemma 2 shows that the candidate approximation is faithful when $C_i$ captures most diffusion mass: the omitted tail controls the unavoidable error, and the learned KL term controls the mismatch inside $C_i$.

\subsection{Spectral Manifold Anchoring}

Diffusion matching transfers local and multi-hop neighborhoods, but a student can still fit candidate distributions while distorting broader community structure. HeatGeo therefore computes spectral coordinates from the teacher graph. Let the symmetric normalized graph Laplacian be
\begin{equation}
    L^T = I - (D^T)^{-1/2}W^T(D^T)^{-1/2}.
\end{equation}
Let $U^T\in\mathbb{R}^{N\times m}$ contain the first $m$ non-trivial eigenvectors of $L^T$, and let $u_i^T$ be the row corresponding to $x_i$. A learnable projection maps student embeddings to this spectral space:
\begin{equation}
    z_i^S = W_{\mathrm{spec}}s_i.
\end{equation}
The spectral loss is
\begin{equation}
    \mathcal{L}_{\mathrm{spec}}
    =
    \frac{1}{|B|}
    \sum_{i\in B}
    \left\|
    z_i^S-u_i^T
    \right\|_2^2.
\end{equation}

\smallskip
\noindent\textbf{Proposition 3 (Spectral anchoring preserves low-frequency topology).}
Fix a mini-batch $B$ and a fixed choice of signs and basis for the teacher Laplacian eigenvectors in $U^T$. If
\begin{equation}
    \mathcal{L}_{\mathrm{spec}}(B)
    =
    \frac{1}{|B|}
    \sum_{i\in B}\left\|z_i^S-u_i^T\right\|_2^2
    \le \eta,
\end{equation}
then the projected student coordinates preserve teacher spectral coordinates in mean square. Moreover, for every pair $i,j\in B$,
\begin{equation}
\begin{aligned}
&\left|\left\|z_i^S-z_j^S\right\|_2
-\left\|u_i^T-u_j^T\right\|_2\right|\\
&\qquad\le
\left\|z_i^S-u_i^T\right\|_2
+\left\|z_j^S-u_j^T\right\|_2,
\end{aligned}
\end{equation}
and therefore
\begin{equation}
\begin{aligned}
&\frac{1}{|B|^2}\sum_{i,j\in B}
\left|\left\|z_i^S-z_j^S\right\|_2
-\left\|u_i^T-u_j^T\right\|_2\right|\\
&\qquad\le 2\sqrt{\eta}.
\end{aligned}
\end{equation}

\paragraph{Implication for HeatGeo.}
The first non-trivial eigenvectors of the normalized graph Laplacian are low-frequency modes of the teacher semantic graph. They vary smoothly over strongly connected regions and separate broader communities, following the graph-manifold principles behind Laplacian Eigenmaps and Diffusion Maps~\citep{belkin2003laplacian,coifman2006diffusionmaps}. The spectral loss therefore anchors the student to the teacher's global topology.

\subsection{Optimization Objectives}

HeatGeo includes a lightweight pointwise anchor for stability:
\begin{equation}
    \mathcal{L}_{\mathrm{anchor}}
    =
    \frac{1}{|B|}
    \sum_{i\in B}
    \left(1-\cos(W_a s_i,t_i)\right),
\end{equation}
where $W_a$ maps student embeddings to the teacher dimension. This term acts as a boundary condition on the embedding map, while the diffusion and spectral objectives carry the main geometric supervision.

The final objective is
\begin{equation}
\begin{aligned}
    \mathcal{L}_{\mathrm{HeatGeo}}
    ={}&
    \lambda_{1}\mathcal{L}_{\mathrm{diff}}
    +
    \lambda_{2}\mathcal{L}_{\mathrm{spec}}
    \\
    &+
    \lambda_{3}\mathcal{L}_{\mathrm{anchor}}.
\end{aligned}
\end{equation}
We use a diffusion-scale curriculum: the first epoch uses one-step diffusion, the second epoch adds two-step diffusion, and later epochs add the full scale set and spectral anchoring. This avoids forcing long-range, smoother structure before the student has learned reliable local neighborhoods.

\smallskip
\noindent\textbf{Corollary 4 (Joint HeatGeo geometry transfer).}
Consider a mini-batch $B$ and active diffusion scales $\mathcal{R}_e$. Assume all loss weights are positive and the candidate sets capture high teacher diffusion mass, i.e., $\rho_{i,r}(C_i)\ge 1-\epsilon$ for all $i\in B$ and $r\in\mathcal{R}_e$. If the weighted HeatGeo objective is small on $B$, then the student simultaneously controls expected candidate-neighborhood mass error, expected full-corpus diffusion error up to the omitted tail $\epsilon$, pairwise distances in the teacher spectral coordinate system, and the pointwise anchor mismatch.

\paragraph{Implication for HeatGeo.}
This result formalizes the role of the three losses. The diffusion loss transfers local and multi-hop retrieval behavior, candidate restriction makes the transfer computationally feasible when captured mass is high, the spectral loss anchors low-frequency manifold topology, and the anchor loss stabilizes the coordinate map. The guarantee is conditional and average-case; it does not claim exact preservation of every full-corpus distance.

\begin{algorithm}[t]
\caption{HeatGeo Training}
\label{alg:heatgeo}
\begin{algorithmic}[1]
\REQUIRE Corpus $\mathcal{D}$, teacher $T$, student $S_{\theta}$, graph size $k$, diffusion scales $\mathcal{R}$
\STATE Cache teacher embeddings $t_i=T(x_i)$ for all $x_i\in\mathcal{D}$
\STATE Build a mutual $k$NN teacher graph $W^T$
\STATE Compute transition matrix $P^T=(D^T)^{-1}W^T$
\STATE Precompute sparse diffusion candidates and targets $p_{i,r}^T$
\STATE Compute spectral coordinates $u_i^T$ from the graph Laplacian
\FOR{each training epoch}
    \STATE Activate diffusion scales according to the curriculum
    \FOR{each mini-batch $B$}
        \STATE Encode anchors and candidate texts with $S_{\theta}$
        \STATE Compute student candidate distributions $p_i^S$
        \STATE Compute $\mathcal{L}_{\mathrm{diff}}$, $\mathcal{L}_{\mathrm{spec}}$, and $\mathcal{L}_{\mathrm{anchor}}$
        \STATE Update $\theta$ by minimizing $\mathcal{L}_{\mathrm{HeatGeo}}$
    \ENDFOR
\ENDFOR
\RETURN Distilled student encoder $S_{\theta}$
\end{algorithmic}
\end{algorithm}
